\section{Falsification controls and scope}
\label{sec:controls}

The preceding results are deductive.  We nevertheless used two finite
controls to expose indexing, boundary, and divisor mistakes during theorem
development.  Their frozen scope and minimum margins appear in
\cref{tab:controls}; the scripts are distributed with the internal artifact
as \texttt{code/verify\_reflection\_rate.py} and
\texttt{code/verify\_least\_period.py}.

\begin{table}[H]
\caption{Finite controls rerun on 25 August 2026.  They are regression and
falsification checks, not substitutes for the proofs.  The least-period
margin compares the exact integer count with the floating-point evaluation
of the analytic lower bound.}
\label{tab:controls}
\centering
\small
\renewcommand{\arraystretch}{1.12}
\begin{tabularx}{\textwidth}{@{}>{\raggedright\arraybackslash}p{0.19\textwidth}
  >{\raggedright\arraybackslash}p{0.41\textwidth}
  >{\raggedright\arraybackslash}p{0.19\textwidth}X@{}}
\toprule
Control & Exhausted family & Minimum recorded margin & Status \\
\midrule
Reflection rate
& 1,093 nonempty symmetric binary adjacency matrices; 5,465 matrix/size cases
& exact integer slack $63$ & PASS \\
Least period
& 25,717 primitive positive-entropy binary matrices of orders $2$--$4$;
514,340 matrix/period pairs
& $2.27201964951$ & PASS \\
\bottomrule
\end{tabularx}
\end{table}


The reflection control exhausts nonempty symmetric $0$--$1$ adjacency
matrices of orders two through four and checks the one-dimensional integer
specialization
\[
 P_{2n-2}q^6\ge C_n^2,
 \qquad P_{2n-2}\le C_{2n-2},
\]
for the recorded sizes.  The exponent $6$ is $2b_1(n)$, since $b_1(n)=3$.
The least-period control computes $F_n=\tr(A^n)$ and exact M\"obius counts as
integers for all primitive positive-entropy binary matrices through order
four.  Eigenvalues enter only the floating-point regression comparison with
\cref{eq:least-bound}; the theorem does not depend on that numerical step.

These controls have three deliberately limited roles.  They can falsify a
proposed universal inequality on the exhausted family, detect off-by-one
errors in the period cell, and preserve a regression target for later edits.
They cannot establish a universal theorem, certify novelty, or validate an
unexamined floating-point implementation.

\begin{example}[Range-two reflection firewall]
\label{ex:range-two}
Consider the one-dimensional reflection-invariant range-two rule
\begin{equation}
 x_i\ne x_{i+2} \qquad(i\in\ZZ).
 \label{eq:range-two-rule}
\end{equation}
Reflecting a block across the site $n$ makes the positions $n-1$ and $n+1$
read the same source position $n-1$.  The reflected pattern therefore has
$x_{n-1}=x_{n+1}$ and violates \cref{eq:range-two-rule}.  Merely sharing a
thicker boundary word does not repair this self-copying obstruction.  Thus
\cref{thm:fixed-period} is a nearest-neighbour result; it must not be quoted
as a general finite-range theorem.
\end{example}

\paragraph{Further limitations.}
The error constants in \cref{eq:entropy-sandwich} depend on the dimension,
and no dimension-uniform rate is claimed.  The shell construction supplies a
lower bound rather than an exact formula for $P_{2n-2}$.  In the edge-shift
setting, the spectral certificate is restricted to primitive matrices with
$\lambda>1$ and may fail to detect small realized periods.  The algorithm in
\cref{thm:least-period} is effective from algebraic spectral data but is not
optimized for numerical conditioning or for the smallest conductor.

\paragraph{Relation to existing mechanisms.}
Bui's reflected periodic extension and explicit entropy estimates
\cite{Bui2026} are the source engine for \cref{sec:reflection}; the residual
step here is the thickness-two common-fibre count with a recoverability
argument and its specified-period normalization.  In
\cref{sec:least-period}, the trace identity, Perron--Frobenius theory, and
M\"obius inversion are standard \cite{LindMarcus2021}.  The residual is the
closed certificate \cref{eq:least-bound} and the rational tail procedure.
No statement in this section upgrades this ownership accounting into a
priority claim.

